\documentclass[11pt]{article}
\usepackage[margin=0.88in]{geometry}
\usepackage[T1]{fontenc}
\usepackage{amsmath,amssymb,hyperref,enumitem}
\hypersetup{hidelinks}
\setlength\parindent{0pt}
\setlength\parskip{0.45em}
\newcommand{\code}[1]{\texttt{\detokenize{#1}}}
\title{The decorated four-to-two crossing, before contraction}
\author{Contextual Composition: ambient pseudomonad development, issue 105}
\date{20 November 2026}
\begin{document}\maketitle
\section*{Audit and interface}
The base is an ordinary category. Grading is split. The objects of the ambient bicategory are strict category-valued presheaves; its 1-cells are strong transformations, not only strict transformations, and its 2-cells are arbitrary modifications. The particular 1-cell inside the square is the native multiplication acting on such an object. No ambient cartesian reflection by multiplication, strict replacement of the comparison, representability, or Kleisli reconstruction is assumed. The assumptions on multiplication are already those of \code{ContextadTotal.Native}; the pre-final comparison needs neither native triangle nor native pentagon.

The test is pointwise on the flat chart of four grades with its terminal presheaf decoration. Equality \emph{after endpoint transport} means a commuting square with equality casts at both ends. It is stronger than being parallel up to arbitrary isomorphism. The component is the inverse multiplication-naturality square at multiplication itself, observed before any final contraction. A unit, associator, or pentagon of horizontal modifications is a different type of output.

\section*{Result}
For every eligible presheaf $X$ and native multiplication $m$, \code{AmbientFourToTwo.exchange} inhabits the previously declared \code{AmbientExchangeDemand.Law X m}. The theorem is not a biconditional with unverified tests. In the flat chart of four inputs to two outputs, at every $x$, the actual inverse square reads the prescribed native disjoint cell. \code{AmbientFourToTwo.row} consequently returns the head-grade transport, a terminal-grade arrow aligned with the measured last grade that factors through the originally fixed right cone with the specified projected base arrow, and a pure terminal $X$-payload. Here pure means an equality transport, not an arbitrary invertible fibre arrow. \code{AmbientFourToTwo.whisker} gives the resulting reading after multiplication, as a consequence, not as evidence for the pre-final assertion.

The passage back to the original interface is deliberate. In the nested chart, the arrows of a twice-applied action have intermediate objects written as iterated one-step powers; the flat chart and the official interchange need not present the same indices syntactically. The proof first reads an arbitrary modification at its prescribed frames, then specializes to the actual square. A predicate indexed by the modification preserves the complete row during coordinate changes. The already established checklist biconditional is used only at the end, with its target fixed to the original law and all of its tests supplied. No final multiplication is cancelled.

\section*{Proof route and fixed witnesses}
Let $f,g:2\to1$ be native routes; assume $g$ takes cartesian grade arrows to cartesian grade arrows. Only at the end set $f=g=\mu$. Write $D_n$ for the flat-to-packed chart, $I$ for splitting $4=2+2$, $J$ for staging $2=1+1$, and $V=D_1$. At a fixed flat input $x$, let $d_x$ be the decorated component of the prescribed disjoint cell. Let $c$ be the last decorated grade of $d_x$, not its projection to the unadorned grade. The following three arrows retain the terminal $X$-decoration:
\[
 u=\widehat{\operatorname{move}}_f(x),\qquad
 r=\widehat{\operatorname{snoc}}_f(\operatorname{inside}_2(g)(x)),\qquad
 \ell=\operatorname{raw}(g)(u).
\]
Here hats are the Grothendieck chosen lifts of the displayed grade arrows; their $X$-payloads are identities. The cast that makes $c;r$ composable is fixed by the serial route frame. In this paragraph $\sim$ denotes equality after equality transport at the source and target, never conjugacy by arbitrary isomorphisms. The grade of $c$ is the prescribed exchange grade, and $\operatorname{grade}(r),\operatorname{grade}(\ell)$ are the original right and left cones. Thus $\operatorname{grade}(c;r)\sim\operatorname{grade}(\ell)$ by the prescribed factorization. Both composites have pure $X$-payload, the latter because restriction maps a pure payload to its functorial image followed by the naturality transport. Their target equality is the inside route's full terminal frame. Uniqueness of a chosen decorated arrow with these target, grade and purity data yields $c;r\sim\ell$. This computation does not yet say which arrow the weak multiplication square chooses.
\par The distinction between a chosen grade arrow and a chosen \emph{packed} arrow matters in the next step. The one-layer chart sends an arrow whose grade and payload have pure fibre components to an arrow with pure packed fibre component. Factor the two-layer chart using $J$ and apply that result twice, once on the last decorated grade and once on the remaining grade, with the stage's actual head and leaf frames. It follows that $D_2(u)$ and $V(r)$ have pure packed fibre components. Their codomains are fixed by the split leaf frames, and their base maps are the prescribed $f$-leg on the original prefix: on the left by \texttt{moveBase}, on the right by \texttt{snocBase} and the explicit inside frame. Chosen-arrow recognition therefore matches these arrows with the actual input and output tips, respectively. It would not match an arbitrary cartesian arrow or an arbitrary strong map by this test.
\par Put $a=\operatorname{oper}(g)$, $q=\operatorname{weakSquare}(f,a)^{-1}_{\rm ambient}$, using the orientation whose raw form runs from left to right, and let $s$ be its nested decorated leaf at $I x$. The intrinsic equation of the weak comparison is $s;\mathrm{tip}_{\rm out}\sim \mathrm{tot}(a)(\mathrm{tip}_{\rm in})$. The source and target tips just computed, together with $D_2;\mathrm{tot}(a)=\operatorname{raw}(g);V$, turn its right side into $V(\ell)$. But $V(c);V(r)\sim V(\ell)$ by the preceding paragraph, and $V(r)\sim\mathrm{tip}_{\rm out}$. Both $s$ and $V(c)$ project to equality transports: the former by the weak comparison's grade projection, the latter by the prescribed exchange's fixed base and the one-layer chart over the base. Cancellation through the original cartesian output tip now gives $s\sim V(c)$. This cancels a cartesian leg over an equal projected arrow; it neither inverts the leg nor asserts uniqueness of arbitrary parallel invertible arrows.
\par Finally the full flat component is recovered by the two parts of $J$: prefix grade and decorated leaf. The actual prefix grade is pure by the weak comparison's grade law after changing charts; the prescribed prefix grade is pure by construction. The leaf comparison above supplies the other part. The two full functor frames fix sources and targets, so joint faithfulness of the pullback projections identifies the staged arrows. Composing $J$ with the nested packed chart gives the flat chart, which is an isomorphism, hence reflects equality. This argument is before the final multiplication and does not ask that multiplication reflect comparisons.

\section*{Horizontal boundary beyond the crossing}
\code{AmbientFlatInside.flat} changes the actual inside block-cell action from a split chart to a flat one, in all arities, given the original cell reading, the two full functor frames of its ends, and the two input/output chart frames. Its one-step unary and ternary specializations admit the native unit and associativity shapes. These are generator-reading rules; they are not ambient horizontal laws. \code{AmbientFlatWeight.reading} does the corresponding job for the actual \emph{canonical strong-arrow comparison}. The named one-step chart views \code{unitView} and \code{muView} include the strict comparison of the charts, not just endpoint-object bijections.

For any strong arrow $f$, \code{AmbientRightWeight.read} further reads the original folded inside-unit square \code{kEta}, not just the canonical comparison to which it is related. The outer fold, the strict pseudofunctor's composition casts and the inserted zero-block frame are accounted for. \code{AmbientUnitSquare.naturality} gives a general route to the unit modification square: for a native $1\to2$ letter, a unit constraint from that letter followed by multiplication to the identity, two full operation frames and an ambient boundary square reading the letter at $f$, it identifies the ambient unit square formed from the serial multiplication boundary and the identity boundary, with the native constraint reframed at both objects. The first-boundary input is not an assumption of the four-to-two theorem. On the right, the data to insert are the inside unit, \code{innerE} and the checked \code{kEta} reading. The exact named public \code{AmbientHorizontalNaturality.Right} specialization is not a declaration in the sealed project. No conclusion for the outer-unit boundary at a general $Kf$, or for the associator square, is claimed from the single crossing.

\section*{Discrimination}
A single test on grade loses the terminal decoration. Conversely, preserving the payload does not force the last grade to factor through the \emph{specified} cone. The predecessor review produced both separations, including a counterfeit with prescribed endpoints, invertibility and a pure payload. Those tests remain applicable; the present argument never reselects a cone after seeing the competing arrow. The new \code{FixedConeCancellation} control also makes the cancellation step falsifiable: in the one-object category of Boolean endofunctions, the nonidentity involution and identity have the same image under the functor to the terminal category and become equal after a constant map; the constant leg is not cartesian. Thus same base and same composite do not suffice without the original cartesian leg. This does not purport to give a counterexample to the native exchange.
\paragraph{Model controls.} The same Boolean endofunction test also uses the identity projection: its constant map is now cartesian, with the same unequal invertible competitors killed after composition; their projected maps now differ. Thus the base equality in cartesian cancellation is independently load-bearing. The supplemental moving-row replay combines the actual positive checklist (hence pure terminal payload) with different source and target objects on the chosen row; \code{noEqualityTransport} shows that \emph{no} arrow between those objects is an equality transport. Terminal purity must not be read as purity of the whole crossing.

The supplemental ordered-grade control rebuilds eleven hash-pinned predecessor fixture modules. \code{Twins.distinction} checks together: the parallel native counterfeit does not align with the prescribed native arrow, its conjugate in the exact ambient frames aligns back to that counterfeit, and the actual ambient arrow aligns with the prescribed one. The pre-existing control proves the conjugate invertible. The generic \code{Twins.rejects} implication explains why these observations distinguish the two ambient arrows; the supplemental claim itself is reported in its checked conjunctive form. These two model probes are additional to the seventeen official probes.
\section*{What does not follow}
A computation at the specific strong arrow $M_X$ is not the modification-naturality theorem for all strong $f:X\to Y$. Neither is a four-to-two calculation (or its postcomposition with $M_X$) an equality between the complete length-four schedules. Before a horizontal triangle or pentagon conclusion, each prescribed outer/inner generator must be read in its proper chart, the \emph{complete} pastings must match the native words with their casts, and the native schedule equality must be invoked. In particular, the result does not promote an objectwise pseudomonad, an identity-on-objects substitution, or a thin observational quotient in place of the requested ambient structure.
\section*{Evidence and replay boundary}
The parent is merge 230, commit \code{22f967a29c13b5bddea550a3b3ccf4f24b94dfae}. Its 259-file source digest manifest has SHA-256
\begin{center}\small\code{f207c3ea28ed6bb05fc638b676def06076d9e6c37f4ac7c4b10ea48a98c3d941}.\end{center}
The project includes the real \code{audit.py} and the three historically missing checks at their official paths; they are verified and preserved, not projected into this parent as extras. From this source the overlay adds 35 production modules and one official probe, updating the umbrella, theorem index and audit aggregate. The 295-source final manifest SHA-256 is
\begin{center}\small\code{9e93c35a700fa403d4e40465f3f430ab09b3d6ab9eec77576d77202447ade586}.\end{center}
Lean is 4.35.0, revision \code{eb6cafb7290e8bd5f2937522d8d804f58f02bafc}; mathlib is revision \code{56a005855ea17e0b0958f73ec801ae140b0a1679}. Production proofs import no predecessor-review probes. The printed axiom transcript, covering 3,750 ordered unique declarations (3,621 parent and 129 new), contains only \code{propext}, \code{Classical.choice} and \code{Quot.sound}; their respective occurrence counts are 3,531, 3,364 and 3,424. Transcript SHA-256:
\begin{center}\small\code{8c45b8a43c993093f338e5cb9539f5eecea1b952a029b8c8017e2f5c763ea84e}.\end{center}
Each new module was recompiled from a fresh copy of its frozen source, into a separate output tree inheriting only the independently rebuilt parent products and pinned upstream packages. The source digest was checked before and after each replay and against the final assembled project. The sixteen unchanged parent probes were run on identical parent source, then each source digest rechecked at final audit; the final new probe, umbrella, full printed axiom audit and audit verifier all ran on the assembled project. This is a warm-parent author replay, not a claim that an entirely empty Lake build was performed in the author phase. The enclosed clean replay script instead assembles a new tree, verifies the pins, then performs a full project build from source. It deliberately does not import author build products.
\end{document}
